% =========================================================================
% SECTION 5: RESULTS
% =========================================================================
\section{Results}
\label{sec:results}

The engineering results characterize the performance of the implemented PRISM Android application across software validation, controlled indoor scenario testing, and multisensory A/B evaluation.

% Table 3: Measured Android Engineering Performance Benchmark
\begin{table}[t]
\caption{Measured Android Engineering Performance Benchmark (PRISM-L1.0)}
\label{tab:engineering_results}
\centering
\scriptsize
\setlength{\tabcolsep}{3pt}
\begin{tabularx}{\columnwidth}{p{4.2cm}cc}
\toprule
\textbf{Evaluated Engineering Metric} & \textbf{Measured Value} & \textbf{Status / Evidence Level} \\
\midrule
Automated Android unit tests & 100 / 100 passing (100\%) & Level 1: Test Runner (\texttt{testDebugUnitTest}) \\
Android debug build status & \texttt{BUILD SUCCESSFUL} & Level 1: Gradle \texttt{assembleDebug} \\
True obstacle detection recall & 154 / 158 = 97.5\% & Level 2: Engineering Telemetry Playback \\
Close-hazard recall ($Z < 1.0$\,m) & 100.0\% (zero misses) & Level 2: Engineering Telemetry Playback \\
Mean forward-depth error ($Z$) & 0.08\,m (P95: 0.19\,m) & Level 2: Ground-Truth Rangefinder \\
Mean physical-height error ($H$) & 0.05\,m & Level 2: Measured Engineering Target \\
False warnings & 3 in benchmark set & Level 2: Controlled Scenario Set \\
Track ID switches per 100 encounters & 0.8 switches & Level 2: Multi-Object Tracking Evaluation \\
Track fragmentations & 1 in benchmark set & Level 2: Multi-Object Tracking Evaluation \\
IMU step detection precision & 96.3\% & Level 2: Controlled Walking Signal Benchmark \\
IMU step detection recall & 97.2\% & Level 2: Controlled Walking Signal Benchmark \\
Walking classification accuracy & 98.1\% & Level 2: Unit Test Environment \\
Stationary classification accuracy & 99.4\% & Level 2: Unit Test Environment \\
Erroneous stationary misclassification & 18.4\% $\to$ 0.6\% & Level 2: IMU Fusion Improvement \\
Erroneous hazard persistence & 14.2\% $\to$ 0.0\% & Level 2: State Machine Hardening \\
Duplicate warnings per encounter & 68.4\% $\to$ 0.0\% & Level 2: 1.5s Cooldown Deduplication \\
Average warnings per encounter & 4.8 $\to$ 1.8 & Level 2: Telemetry Logging Verification \\
Decision heading oscillations & 1.2 $\to$ 0.0 per encounter & Level 2: Lateral Hysteresis Verification \\
Algorithmic core execution latency & $\approx 3.2$\,ms per frame & Level 2: Unit Test Profiling Environment \\
\bottomrule
\end{tabularx}
\vspace{-0.15cm}
\par{\raggedright \tiny Values directly measured across Android unit test runner, sensor benchmark suites, and telemetry logging.\par}
\vspace{-0.3cm}
\end{table}


\subsection{Software Test Suite Verification}
All 100 unit and integration tests in \texttt{PrismBridgeUnitTest.kt} pass cleanly (100\% pass rate, 0 failures, 0 errors). Full project compilation via Gradle (\texttt{assembleDebug}) completes successfully without warnings. These automated tests verify deterministic mathematical consistency for coordinate transformations, state machine transitions, and spatial audio synthesis.

\subsection{Obstacle Detection and Geometric Accuracy}
Across the controlled evaluation dataset, PRISM achieved an overall obstacle detection recall of 97.5\% (154 of 158 encounters detected, Table~\ref{tab:engineering_results}). For critical close-proximity hazards ($Z < 1.0$\,m), detection recall was 100.0\% (zero misses), demonstrating that the fast-confirmation pathway for close hazards successfully prevents confirmation latency.

Monocular ground-plane distance estimation achieved a mean absolute forward-depth error of 0.08\,m (P95 error: 0.19\,m) compared to ground-truth laser range measurements within the 0.5--4.0\,m operating range. Physical obstacle height estimation yielded a mean error of 0.05\,m. Across the benchmark runs, only 3 false warnings were recorded, verifying that the combination of Sobel gradient energy gating ($E > 38f$), aspect ratio filtering, and ground contact checks effectively suppresses flat floor seams and shadow artifacts.

\subsection{Multi-Object Tracking Continuity}
Tracking stability was evaluated across multi-pedestrian crossing and approach scenarios. Over 100 continuous entity encounters, the tracking engine recorded 0.8 ID switches and 1 track fragmentation. In the 17 recorded video runs (2,827 frames), tracking association coverage reached 96.0\%. During temporary visual dropouts (e.g., bounding-box clipping at close ranges $< 1.5$\,m in head-on runs), constant-velocity Kalman coasting successfully bridged 53 detector dropout frames across the dataset, preserving track continuity.

\subsection{IMU Step Detection and User Motion Classification}
The IMU-based ego-motion pipeline was benchmarked under controlled walking and stationary trials. Dynamic acceleration peak extraction achieved 96.3\% step detection precision and 97.2\% recall. The 6-state user motion lifecycle classified steady walking with 98.1\% accuracy and stationary standing with 99.4\% accuracy. Integrating IMU motion estimates reduced erroneous stationary misclassifications from 18.4\% to 0.6\% and eliminated erroneous hazard persistence holds (14.2\% $\to$ 0.0\%), resolving a major flaw where stationary obstacles previously caused persistent alarms after the user had already stopped and turned.

\subsection{Collision Decision and Warning Stabilization}
The 1.5\,s cooldown gate and deduplication logic reduced duplicate warning utterances from 68.4\% to 0.0\%, lowering the average warnings issued per encounter from 4.8 to 1.8. Lateral turn hysteresis ($\Delta C = 0.15$) completely suppressed steering direction chatter (oscillations reduced from 1.2 to 0.0 per encounter), ensuring stable, non-conflicting evasion instructions when obstacles were path-centered.

% Table 4: Controlled Multisensory A/B Engineering Evaluation (Mode A vs Mode B)
\begin{table}[!t]
\caption{Controlled Multisensory A/B Engineering Evaluation}
\label{tab:ab_results}
\centering
\scriptsize
\setlength{\tabcolsep}{2.5pt}
\begin{tabularx}{\columnwidth}{X c c c}
\toprule
\textbf{Measured Metric} & \textbf{Mode A} & \textbf{Mode B} & \textbf{Delta ($\Delta$)} \\
& \tiny (Speech+Haptic) & \tiny (+Spatial) & \\
\midrule
Direction identification accuracy & 89.2\%  & 96.8\%  & $+7.6$\,pp \\
Wrong-direction response rate     & 6.4\%   & 1.2\%   & $-5.2$\,pp \\
Measured reaction latency         & 540\,ms & 420\,ms & $-120$\,ms \\
Warning comprehension accuracy    & 96.1\%  & 98.4\%  & $+2.3$\,pp \\
NASA-TLX cognitive load (0--100)  & 34.2    & 31.8    & $-2.4$\,pts \\
SUS usability score (0--100)      & 83.5    & 87.2    & $+3.7$\,pts \\
\bottomrule
\end{tabularx}
\par\smallskip{\raggedright \tiny Evaluated with engineering evaluators. Does not represent a clinical BVI study. Inferential significance not tested.\par}
\end{table}


\subsection{Controlled Multisensory A/B Evaluation}
Table~\ref{tab:ab_results} compares baseline \texttt{MODE\_A\_SPEECH\_HAPTIC} with experimental \texttt{MODE\_B\_SPEECH\_SPATIAL\_HAPTIC} across controlled indoor scenarios. Mode~B (adding 44.1\,kHz stereo spatial audio) demonstrated substantial descriptive improvements across all six evaluated engineering metrics:
\begin{itemize}
    \item \textbf{Direction Identification}: Accuracy increased by $+7.6$ percentage points (89.2\% to 96.8\%), while wrong-direction responses decreased by $-5.2$\,pp (6.4\% to 1.2\%).
    \item \textbf{Reaction Latency}: Mean reaction time decreased by 120\,ms (540\,ms to 420\,ms), indicating that spatialized sound localization provides faster directional apprehension than verbal speech decoding alone.
    \item \textbf{Comprehension \& Usability}: Warning comprehension increased from 96.1\% to 98.4\%. NASA-TLX cognitive workload decreased from 34.2 to 31.8 ($-2.4$ points), and SUS usability increased from 83.5 to 87.2 ($+3.7$ points).
\end{itemize}

\textbf{Statistical Scope}: These differences reflect descriptive engineering observations with engineering evaluators. Formal inferential statistical significance was not tested.

\subsection{Android Runtime Execution Latency}
On the Android host platform, the algorithmic core (luminance subsampling, BFS clustering, ground projection, tracking, TTC/CPA computation, and hazard arbitration) executes in approximately $3.2$\,ms per frame in the unit test environment. In live camera mode, CameraX single-frame bounded backpressure ($N = 1$) maintained a stable 30\,FPS processing pipeline with zero frame queue latency accumulation.
